\clearpage
\begin{summarybox}
  \textbf{\textcolor{CSummary}{一句话核心}}

  \textbf{焦半径不必算斜线：转到准线，沿开口方向的坐标距离加上焦点到顶点的距离。}
  \[
    \boxed{PF=t+\frac p2},\qquad t\ge0.
  \]
  对 $y^2=2px$，$t=x_0$，所以 \textbf{$PF=x_0+\frac p2$}。

  \medskip
  \textbf{\textcolor{CSummary}{使用条件}}
  \begin{itemize}[leftmargin=*, itemsep=0.35em, topsep=0.35em]
    \item 抛物线为标准式，$p>0$；\textbf{点必须在该抛物线上}。
    \item 向右 $t=x_0$；向左 $t=-x_0$；向上 $t=y_0$；向下 $t=-y_0$。
  \end{itemize}

  \medskip
  \textbf{\textcolor{CSummary}{三种证明，按需求选}}
  \begin{description}[style=nextline, leftmargin=2.8em, labelsep=0.5em, itemsep=0.35em]
    \item[\textbf{最快：}] $PF=d(P,l)=x_0-(-p/2)=x_0+p/2$（向右）。
    \item[\textbf{门槛最低：}] 水平走两段：准线 $\to y$ 轴为 $p/2$，$y$ 轴 $\to P$ 为 $x_0$。
    \item[\textbf{常规：}] $PF=\sqrt{(x_0-p/2)^2+y_0^2}$，代 $y_0^2=2px_0$ 凑完全平方。
  \end{description}

  \medskip
  \textbf{\textcolor{CSummary}{绝不混淆的三个量}}
  \[
    \underbrace{OF}_{\text{焦点到顶点}}=\frac p2,\qquad
    \underbrace{d(F,l)}_{\text{焦准距}}=p,\qquad
    \underbrace{\text{通径长}}_{\text{通过焦点、垂直对称轴的弦}}=2p.
  \]
  半通径是 $p$，不是 $p/2$。

  \medskip
  \textbf{\textcolor{CSummary}{三秒自检}}
  $P=O$ 时 $PF=p/2$（最短）；$x_0\ge0$（向右）；
  点不在曲线上不能直接套用。
\end{summarybox}
